\chapter{Kurzschlussberechnung}
\label{chap:08-kurzschlussberechnung}

%\chapter{Dreipolige Kurzschlussberechnung}
%\label{chap:kurzschluss}

\section{Ziel und Einordnung}

Ein dreipoliger Kurzschluss verbindet die drei Phasen an einer Fehlerstelle.
Bei ideal symmetrischen Netzbedingungen ist der Fehler symmetrisch.
Die Berechnung kann dann mit einem einphasigen Ersatzkreis erfolgen.

Der dreipolige Kurzschluss führt häufig zu hohen Strömen und ist für die
Bemessung von Schaltgeräten, Sammelschienen, Leitern und Schutzsystemen
von Bedeutung. Er ist jedoch nicht in jedem Netz und für jede Fehlerart
der Fehler mit dem höchsten Strom.

\section{Ersatzschaltbild und Anfangskurzschlussstrom}

Das Netz wird aus Sicht der Fehlerstelle durch eine Ersatzspannung und
eine Ersatzimpedanz beschrieben. Für den symmetrischen Drehstromfall
ist die Strangspannung
\[
  U_{\mathrm{ph}}=\frac{U_{LL}}{\sqrt{3}}.
\]

Der anfängliche symmetrische Kurzschlussstrom kann in einem vereinfachten
Ersatzmodell durch
\[
  I_k''=\frac{U_{\mathrm{ph}}}{|Z_k|}
  =\frac{U_{LL}}{\sqrt{3}|Z_k|}
\]
bestimmt werden.

Das Symbol \(I_k''\) wird in der Kurzschlussberechnung häufig für den
anfänglichen symmetrischen Kurzschlussstrom verwendet. Bei normgerechten
Berechnungen sind die jeweils geltenden Definitionen und Korrekturfaktoren
zu berücksichtigen.

\section{Kurzschlussleistung}

Die dreiphasige Kurzschlussleistung wird häufig definiert durch
\[
  S_k''=\sqrt{3}\,U_{LL}I_k''.
\]

Daraus folgt
\[
  I_k''=\frac{S_k''}{\sqrt{3}\,U_{LL}}.
\]

Unter der vereinfachenden Annahme eines symmetrischen Netzes und einer
geeigneten Ersatzimpedanz kann der Betrag der Netzimpedanz abgeschätzt
werden als
\[
  |Z_k|\approx\frac{U_{LL}^2}{S_k''}.
\]

Die Kurzschlussleistung ist eine Rechengröße, die den Kurzschlussstrom
mit der Bezugsspannung verknüpft. Sie ist nicht mit der während des
gesamten Fehlerverlaufs konstant abgegebenen Wirkleistung gleichzusetzen.

\section{Beispiel: Kurzschlussleistung auf Mittelspannungsebene}

Gegeben seien
\[
  U_{LL}=20\,\mathrm{kV},\qquad S_k''=250\,\mathrm{MVA}.
\]

Der anfängliche symmetrische Kurzschlussstrom beträgt näherungsweise
\[
  I_k''=\frac{250\,\mathrm{MVA}}
  {\sqrt{3}\cdot20\,\mathrm{kV}}
  \approx7{,}22\,\mathrm{kA}.
\]

Der Betrag der entsprechenden Ersatzimpedanz ist
\[
  |Z_k|\approx\frac{(20\,\mathrm{kV})^2}{250\,\mathrm{MVA}}
  =1{,}6\,\Omega.
\]

Die Impedanz ist damit nicht vollständig beschrieben: Für eine genaue
Phasorrechnung werden auch der Wirk- und Blindanteil benötigt.

\section{Wirk- und Blindanteil}

Die Kurzschlussimpedanz wird als
\[
  Z_k=R_k+jX_k
\]
geschrieben. Ihr Betrag ist
\[
  |Z_k|=\sqrt{R_k^2+X_k^2}.
\]

Das Verhältnis \(X/R\) beeinflusst den zeitlichen Verlauf des Kurzschlussstroms,
insbesondere den Gleichstromanteil und damit die Stoßstrombeanspruchung.

Wird ausschließlich der Betrag der Impedanz verwendet, kann der anfängliche
symmetrische Wechselstromanteil abgeschätzt werden. Für die Berechnung
des Stoßkurzschlussstroms oder der Ausschaltbeanspruchung ist die
Impedanzphase beziehungsweise das Verhältnis \(X/R\) jedoch wichtig.

\section{Beitrag verschiedener Einspeiser}

In einem realen Netz können Kurzschlussströme aus mehreren Quellen stammen:
\begin{itemize}
  \item dem übergeordneten Versorgungsnetz;
  \item Synchronmaschinen;
  \item Asynchronmaschinen, abhängig von Betriebszustand und Fehlerdauer;
  \item leistungselektronischen Einspeisern mit regelungsabhängigem
        Fehlerstromverhalten.
\end{itemize}

Die Beiträge dürfen nicht unabhängig von ihrer Phasenlage und den
Netzimpedanzen als bloße Beträge addiert werden. Außerdem unterscheiden
sich ihre zeitlichen Verläufe. Für die Bemessung von Schaltgeräten müssen
Anfangskurzschlussstrom, Stoßstrom, Ausschaltstrom und thermische
Beanspruchung entsprechend dem Berechnungsziel getrennt betrachtet werden.

Bei umrichterbasierten Quellen kann ein einfaches Modell als ideale
Spannungsquelle hinter einer festen Impedanz unzureichend sein. Das
Fehlerstromverhalten hängt von Strombegrenzung, Regelung und Schutzstrategie ab.

\section{Thermische und dynamische Beanspruchung}

Der Kurzschlussstrom erzeugt Joule-Verluste, die bei ausreichender Fehlerdauer
zu einer erheblichen Erwärmung führen können. Die elektrodynamischen Kräfte
zwischen stromführenden Leitern steigen näherungsweise mit dem Quadrat
des momentanen Stroms.

Für die Auswahl von Betriebsmitteln müssen deshalb mindestens unterschieden
werden:
\begin{itemize}
  \item der anfängliche symmetrische Kurzschlussstrom;
  \item der Stoßkurzschlussstrom;
  \item der während der Ausschaltzeit wirksame Kurzschlussstrom;
  \item die thermische Beanspruchung über die Fehlerdauer;
  \item das Ausschaltvermögen und die Kurzzeitstromfestigkeit des Betriebsmittels.
\end{itemize}

Die konkrete Bemessung richtet sich nach dem jeweiligen Betriebsmittel,
der Netzkonfiguration und den anwendbaren Normen.

\section{Rechenablauf}

Ein nachvollziehbarer Rechenablauf umfasst:
\begin{enumerate}
  \item Fehlerstelle und betrachteten Betriebszustand festlegen.
  \item Nennspannung und erforderliche Kurzschlusskenngröße bestimmen.
  \item Netz-, Transformator- und Leitungsimpedanzen zusammentragen.
  \item Impedanzen auf eine gemeinsame Spannungsebene beziehen.
  \item Wirk- und Blindanteile komplex addieren.
  \item Den symmetrischen Kurzschlussstrom berechnen.
  \item Bei Bedarf Stoßstrom, Ausschaltstrom und thermische Beanspruchung
        nach dem geeigneten Verfahren bestimmen.
  \item Ergebnis auf Einheiten, Plausibilität und Vollständigkeit prüfen.
\end{enumerate}

\section{Grenzen der vereinfachten Berechnung}

Die einfache Gleichung
\[
  I_k''=\frac{U_{LL}}{\sqrt{3}|Z_k|}
\]
ist ein Ersatzmodell. Sie berücksichtigt nicht automatisch alle
Einflussgrößen einer vollständigen normativen Kurzschlussberechnung.

Je nach Netz und Berechnungsziel können erforderlich sein:
\begin{itemize}
  \item Spannungsfaktoren und Betriebsmittelkorrekturen;
  \item maximale und minimale Netzkonfigurationen;
  \item Widerstandsänderung durch Erwärmung;
  \item Motor- und Generatorbeiträge;
  \item Umrichtermodelle;
  \item die zeitliche Entwicklung des Fehlerstroms;
  \item die tatsächlich vorhandene Fehlerimpedanz.
\end{itemize}

Für Schutzkoordination ist außerdem nicht nur der maximale Kurzschlussstrom
relevant. Der minimale Fehlerstrom kann entscheidend dafür sein, ob ein
Schutzgerät innerhalb der erforderlichen Zeit auslöst.

\section{Übungsaufgaben}

\begin{enumerate}
  \item Berechnen Sie den anfänglichen symmetrischen Kurzschlussstrom
        bei \(U_{LL}=10\,\mathrm{kV}\) und \(S_k''=200\,\mathrm{MVA}\).
  \item Eine Netzimpedanz beträgt \(Z_k=0{,}2+j1{,}0\,\Omega\).
        Bestimmen Sie ihren Betrag und das Verhältnis \(X/R\).
  \item Warum reicht der Betrag der Kurzschlussimpedanz nicht aus,
        um alle Kenngrößen des Kurzschlussverlaufs zu bestimmen?
  \item Erklären Sie den Unterschied zwischen Kurzschlussstrom und
        Kurzschlussleistung.
  \item Warum müssen für die Schutzkoordination sowohl maximale als
        auch minimale Fehlerströme betrachtet werden?
\end{enumerate}
